\chapter{Weakly Interacting Bose Gase}

\section{The Weak Interaction Assumption}

We take a box with periodic boundary conditions and volume $V=L^3$.\TODO{Add a figure of the box.} Inside we consider a system of $N$ bosons with mass $m$ and a two-body interaction potential $U(\mathbf{r})$. There are systems of bosons that can be described by this model, such as liquid helium-4 (which actually is a strongly interacting system) and ultracold atomic gases. The Hamiltonian of the system is given by
\[
    \hat{H} = \sum_{i= 1}^N \frac{\hat{\mathbf{p}}_i^2}{2m} + \sum_{i<j} U(\hat{\mathbf{r}}_i - \hat{\mathbf{r}}_j),
\]

[...]

... six bullet points... now how to replace the interaction term with a simpler one.

We replace \(U(\mathbf{q})\)  with \(U(0)\), since if the potntial is sufficiantly smooth, \textbf{short-ranged} and \textbf{monotonic}, the Fourier transform of the potential is approximately constant for small momenta. Then the main contribution to the interaction term comes from the zero-momentum component of the potential. In the context of two body interactions, wan der Waals interactions are a good example of a short-ranged potential. IN the Lennard Jones potential, the attractive part of the potential is proportional to \(1/r^6\) and the repulsive part is proportional to \(1/r^{12}\). We have also to remember dipole interactions, which are proportional to \(1/r^3\) (which can be tricky since it is a marginal case). Since we are dealing with atomes, Wan der Waals interactions are effectively short-ranged, since the typical distance between atoms is much larger than the atomic radius. Since pions are the mediators of the strong force and they are massive, the strong force is also effectively short-ranged.

We can now approximate the operators \(\hat{a}_0\) and \(\hat{a}_0^\dagger\) by their expectation values, which are \(\sqrt{N_0}\) and \(\sqrt{N_0}\), respectively... hamiltonian...

The terms which cancels out are the ones that contain \(\hat{a}_0\) and \(\hat{a}_0^\dagger\) in the interaction term: if it wasn't the case they could have contributed to renormalization of the free term of the hamiltonian.

Now we have only creation operators, not number operators, so we can employ the Bogoliubov transformation to diagonalize the Hamiltonian.

\subsection{Bogoliubov Transformations}

We need to condtruct an hamiltonian which in the end has a number operator in the form \(\hat{a}_\mathbf{k}^\dagger \hat{a}_\mathbf{k}\). We can do this by performing a Bogoliubov transformation, which is a linear transformation of the creation and annihilation operators.

We start with an ansatz for the Bogoliubov transformation:

[...]

commutation relations... condition on coefficents and transformation matrix (to invert ansatz)...

diagonalization of the Hamiltonian... conditions on arbitrary coefficients... final form of the Hamiltonian...

\subsubsection{Ground State}

[...]